\documentclass[10pt,a4paper]{amsart}
\usepackage[margin=19mm]{geometry}
\usepackage{amsmath,amssymb,mathtools}
\usepackage[scr=rsfs,cal=euler]{mathalfa}
\usepackage{xfrac}
\usepackage[unicode,hidelinks,bookmarks=false]{hyperref}
\newcommand{\ie}{i.e.\ }
\newcommand{\cf}{cf.\ }
\newcommand{\resp}{respectively\ }
\newcommand{\Subsection}{\subsection}
\providecommand{\nobreakdash}{\nobreak\hbox{-}}
\setlength{\emergencystretch}{2em}
\multlinegap0pt
\usepackage[shortlabels]{enumitem}
\usepackage{amssymb}
\usepackage{mathtools}
\usepackage{tikz}
\usepackage{xcolor}
\newtheorem{corollary}{Corollary}
\newtheorem{lemma}{Lemma}
\newtheorem{theorem}{Theorem}
\usepackage{cleveref}
\crefname{corollary}{Corollary}{Corollarys}
\crefname{lemma}{Lemma}{Lemmas}
\crefname{theorem}{Theorem}{Theorems}
\DeclareMathOperator{\Ho}{H}
\newcommand{\bbd}{\mathbb{D}}
\newcommand{\calfc}{\mathcal{FC}}
\newcommand{\calfm}{\mathcal{F\!M}}
\title{L-fuzzy simplicial homology}
\author{Javier Perera-Lago, Alvaro Torras-Casas, Rocio Gonzalez-Diaz}
\date{Source: arXiv:2604.08166v1 (CC BY 4.0)}
\begin{document}
\maketitle

\noindent\emph{Source and licence.} L-fuzzy simplicial homology. Version arXiv:2604.08166v1 (CC BY 4.0), distributed by arXiv under the Creative Commons Attribution 4.0 International licence, http://creativecommons.org/licenses/by/4.0/, as recorded in the arXiv licence metadata for this exact version and shown on the abstract page https://arxiv.org/abs/2604.08166v1. Full licence notice: \texttt{licenses/arxiv-2604.08166-CC-BY-4.0.txt}.\par\smallskip

\noindent\textit{Editorial scope.} Counted unit: the article's theorem thm:hdl\_properties (source0.tex lines 1282--1306). The article's complete original proof of it is delivered with it (source0.tex lines 1308--1370, one \texttt{proof} environment of 63 lines). No mathematics is written, repaired or re-derived: the statement and the proof are the article's own lines, in the article's own notation, and no retained line is substituted or re-flowed.  Retained uncounted as the source-local context the proof needs: lines 880--889; lines 1177--1180; lines 1233--1241.  Nothing was omitted from any retained span, so the package carries no omission record.  No retained line contains a citation, so the wrapper carries no bibliography and no bibliography entry.  No retained line contains a figure environment, an image-inclusion command or drawing code, so no image is delivered and the package declares no asset dependency.  Editorial formatting only, in addition: an AMS \texttt{amsart} preamble that restates the article's own declarations and macros used by the retained lines, namely the environments \texttt{corollary}, \texttt{lemma}, \texttt{theorem} on separate counters, together with the standard packages those lines need.  The retained environments print in this document as Corollary 1, Lemma 1, Theorem 1 in order of appearance, whereas the article prints the same items with the numbering of its own full document.  The complete native source of the article as distributed in the arXiv e-print for this exact version is bundled unchanged as \texttt{sources/198140/source0.tex}.
\begin{corollary}
\label{cor:eta_value}
Let \(\Delta\) be a simplicial complex, let \((L,\leq)\) be a CDL such that \(0 \in L\) is meet-prime and let \(\mu \in \calfc(\Delta,L)\) such that  \(\mu^* = \Delta\). If \(\Ho_d\) is the \(d\)-homology \(\bbd\)-module of \(\Delta\), then \(\eta_d \in \calfm(\Ho_d,L)\) and, for each \([h] \in \Ho_d\),
\[
\eta_d([h]) =
\bigvee \{ \zeta_d(z) \mid z \in [h] \}
=\bigvee \{ \kappa_d(z) \mid z \in [h]\}.
\]
In particular, $\eta_d^*=\Ho_d$.
\end{corollary}

\begin{lemma}
\label{lem:shi_welldefined}
Let \([h] \in \Ho_d\) and let \(z \in [h]\). Then, \(S(h,I)\) is solvable if and only if \(S(z,I)\) is solvable.
\end{lemma}

\begin{lemma}
\label{lem:threshold_system}
Let \(\Delta\) be a simplicial complex, let \(\mu \in \calfc(\Delta,L)\), and let \([h] \in \Ho_d\). For each \(\ell \in L\),
define the index subset
\(I(\ell) = \big\{i \in \{1, \ldots, n_d\}\mid \sigma^{d}_{i} \in \mu^{\not\geq \ell}\big\}.\)
Then, the vector 
\(\left( \begin{array}{c} \upsilon \\ \tau \end{array} \right) \in \bbd^{n_U+n_T}\)
is a solution for the system \(S(h, I(\ell))\) if and only if the \(d\)-cycle \(z \in [h]\) with \(z \approx h+(\upsilon,\tau)\) satisfies that \(\kappa_d(z)\geq \ell\).
\end{lemma}

\begin{theorem}
\label{thm:hdl_properties}
Let \(\Delta\) be a simplicial complex and \(\Ho_d\) its \(d\)-homology \(\bbd\)-module. Let \(\mu \in \calfc(\Delta,L)\) be an \(L\)-fuzzy subcomplex of \(\Delta\) and \(\eta_d \in \calfm(\Ho_d,L)\) the \(L\)-fuzzy \(d\)-homology of \(\mu\). For each \(\ell \in L\),
define the crisp subset of \(\Ho_d\):
\[
\Ho_d(\ell) = \big\{[h] \in \Ho_d \mid S(h,I(\ell)) \text{ is solvable} \big\}.
\]
We claim:
\begin{enumerate}[label=(\roman*)]
\item \label{item:hdl_submodule} The set \(\Ho_d(\ell)\) is a crisp submodule of \(\Ho_d\).
\item \label{item:hdl_functor} For any \(\ell_1, \ell_2 \in L\) with \(\ell_1 \leq \ell_2\), we have \(\Ho_d(\ell_1) \supseteq \Ho_d(\ell_2)\).
\item \label{item:hdl_join} For any subset \(S \in L\), we have \(\Ho_d(\bigvee S)\subseteq \bigcap_{\ell \in S} \Ho_d(\ell)\).
\item \label{item:hdl_meet} For any subset \(S \in L\), we have \(\Ho_d(\bigwedge S)\supseteq \bigcap_{\ell \in S} \Ho_d(\ell)\).
\item \label{item:hdl_eta} For every \([h] \in \Ho_d\),
\[
\eta_d([h]) = \bigvee \Big\{\, \ell \in L(\kappa_d) \, \Big| \, [h] \in \Ho_d(\ell) \,\Big\}.
\]
\item \label{item:hdl_contain} For any \(\ell \in L\), we have \(\Ho_d(\ell) \subseteq \eta_d^{\geq \ell}\).
\item \label{item:hdl_cut} For any \(\ell \in L\),
\[
\eta_d^{\geq \ell}
= \bigcup_{\substack{S \subseteq L(\kappa_d) \\ \bigvee S \geq \ell}} \bigcap_{s \in S} \Ho_d(s).
\]
\end{enumerate}
\end{theorem}

\begin{proof}
First note that \(\Ho_d(\ell)\) is well-defined, since \Cref{lem:shi_welldefined} ensures that the solvability of \(S(h,I(\ell))\) is independent of the chosen representative of \([h]\).
\begin{enumerate}[label=(\roman*)]
\item The zero class \([0] \in \Ho_d\) belongs to \(\Ho_d(\ell)\) because \(S(0, I(\ell))\) is homogeneous and trivially solvable.
Let \([h_1], [h_2] \in \Ho_d(\ell)\). Then \(S(h_1, I(\ell))\) and \(S(h_2, I(\ell))\) have solutions 
\(\left(\begin{array}{c}\upsilon_1 \\ \tau_1\end{array}\right)\)
and 
\(\left(\begin{array}{c}\upsilon_2 \\ \tau_2\end{array}\right)\) 
respectively. Hence:
\[
(U_{d} \mid T_{d} A_d)_{I(\ell)}
\left(\begin{array}{c}\upsilon_1 \\\tau_1\end{array}\right)= -(h_1)^{\Delta}_{I(\ell)}, 
\quad
(U_{d} \mid T_{d} A_d)_{I(\ell)}
\left(\begin{array}{c}\upsilon_2 \\\tau_2\end{array}\right)= -(h_2)^{\Delta}_{I(\ell)}.
\]
Adding these equations yields:
\[
(U_{d} \mid T_{d} A_d)_{I(\ell)}
\left(\begin{array}{c}\upsilon_1 + \upsilon_2 \\\tau_1 + \tau_2\end{array}\right)= -(h_1 + h_2)^{\Delta}_{I(\ell)}.
\]
Thus, \(S(h_1 + h_2, I(\ell))\) is solvable, so \([h_1 + h_2] = [h_1] + [h_2] \in \Ho_d(\ell)\). Now, let \([h] \in \Ho_d(\ell)\) with a solution 
\(\left(\begin{array}{c}\upsilon \\ \tau\end{array}\right)\), and let \(a \in \bbd\). Then:
\[
(U_{d} \mid T_{d} A_d)_{I(\ell)}
\left(\begin{array}{c}a\upsilon \\a\tau\end{array}\right)= a \cdot
(U_{d} \mid T_{d} A_d)_{I(\ell)}
\left(\begin{array}{c}\upsilon \\\tau\end{array}\right)= -(a h)^{\Delta}_{I(\ell)}.
\]
Hence \(S(ah, I(\ell))\) is solvable and therefore \([ah] = a[h] \in \Ho_d(\ell)\).
Consequently, \(\Ho_d(\ell)\) is a submodule of \(\Ho_d\).
\item If \(\ell_1 \leq \ell_2\), then \(I(\ell_1) \subseteq I(\ell_2)\). Thus, for any \([h] \in \Ho_d\), the system \(S(h,I(\ell_2))\) contains all equations of \(S(h,I(\ell_1))\). Therefore, the solvability of \(S(h,I(\ell_2))\) implies the solvability of \(S(h,I(\ell_1))\), and hence \(\Ho_d(\ell_2) \subseteq \Ho_d(\ell_1)\).

\item From \cref{item:hdl_functor} we have \(\Ho_d(\ell) \supseteq \Ho_d \left(\bigvee S\right)\) for all \(\ell \in S\), which implies that \(\Ho_d(\bigvee S) \subseteq \bigcap_{\ell \in S} \Ho_d(\ell)\).

\item From \cref{item:hdl_functor} we have \(\Ho_d(\ell) \subseteq \Ho_d \left(\bigwedge S\right)\) for all \(\ell \in S\), which implies that \(\Ho_d(\bigwedge S) \supseteq \bigcup_{\ell \in S} \Ho_d(\ell)\).

\item\label{item:Lh-Sh} Having fixed \([h] \in \Ho_d\), define the sets:
\[
L_h = \big\{\kappa_d(z) \mid z \in [h]\big\} \subseteq L(\kappa_d), 
\qquad 
S_h = \big\{\, \ell \in L(\kappa_d) \mid [h] \in \Ho_d(\ell) \,\big\} \subseteq L(\kappa_d).
\]
By \Cref{cor:eta_value}, we know that \(\eta_d([h])=\bigvee L_h\). We now show that \(\bigvee L_h = \bigvee S_h\) by proving both inequalities.
If \(\ell \in L_h\), then there exists \(z \in [h]\) such that \(\kappa_d(z) = \ell\). By \Cref{lem:threshold_system}, \(S(h,I(\ell))\) is solvable, hence \([h] \in \Ho_d(\ell)\) and \(\ell \in S_h\). Thus \(L_h \subseteq S_h\), and consequently \(\bigvee L_h \leq \bigvee S_h\). Conversely, if \(\ell \in S_h\), then \([h] \in \Ho_d(\ell)\) and \(S(h,I(\ell))\) is solvable. By \Cref{lem:threshold_system}, there exists \(z \in [h]\) such that \(\kappa_d(z) \geq \ell\), with \(\kappa_d(z) \in L_h\). Therefore, each element of \(S_h\) is bounded by an element of \(L_h\), implying \(\bigvee S_h \leq \bigvee L_h\). Hence \(\eta_d([h])=\bigvee L_h = \bigvee S_h\).
\item Let \([h] \in \Ho_d(\ell)\).
Then $\ell \in S_h$; where $S_h$ is defined as in \cref{item:Lh-Sh}. 
Further, by \cref{item:hdl_eta}, we have \(\eta_d([h])=\bigvee S_{h}\). Hence \(\ell \leq \eta_d([h])\), which implies that \([h] \in \eta_d^{\geq \ell}\). Therefore \(\Ho_d(\ell) \subseteq \eta_d^{\geq \ell}\).
\item We prove both inclusions. First assume that \([h] \in \eta_d^{\geq \ell}\).
By \cref{item:hdl_eta}, we have \(\eta_d([h])=\bigvee S_h \geq \ell\). Since \([h] \in \Ho_d(s)\) for every \(s \in S_h\), it follows that \([h] \in \bigcap_{s \in S_h} \Ho_d(s)\). Because \(\bigvee S_h \geq \ell\), we obtain
\[
[h] \in 
\bigcup_{\substack{S \subseteq L(\kappa_d) \\ \bigvee S \geq \ell}}
\bigcap_{s \in S} \Ho_d(s).
\]
Conversely, assume that \([h]\) belongs to that set.
Then there exists a subset \(S \subseteq L(\kappa_d)\) such that \([h] \in \Ho_d(s)\) for all \(s \in S\) and \(\bigvee S \geq \ell\).
In particular, $S \subseteq S_h$.
By \cref{item:hdl_eta},
\(\eta_d([h])= \bigvee S_h\geq \bigvee S\geq \ell.\)
Therefore \([h] \in \eta_d^{\geq \ell}\), which completes the proof. \qedhere
\end{enumerate}
\end{proof}

\end{document}
